\subsection{Marco Conceptual}

Los siguientes conceptos t\'ecnicos fundamentan el proyecto. Algunos corresponden a componentes efectivamente utilizados en la aplicaci\'on final y otros a estrategias de optimizaci\'on evaluadas durante el dise\~no del \textit{pipeline}; su presencia en este glosario no implica que todos est\'en expuestos como funcionalidad visible en la build final.

\textit{Albedo.} Mapa de textura que define el color base de un material sin incorporar iluminaci\'on. Es una de las entradas principales del flujo PBR.

\textit{Asset Bundle.} Mecanismo de empaquetado de Unity para distribuir activos en bloques independientes. En este proyecto se consider\'o como alternativa de despliegue, pero no constituye el eje central documentado para la versi\'on final de la aplicaci\'on.

\textit{Baking.} Proceso de precalcular informaci\'on compleja, como detalle geom\'etrico o iluminaci\'on, y almacenarla en texturas para reducir costo computacional en tiempo real.

\textit{BRDF (\textit{Bidirectional Reflectance Distribution Function}).} Funci\'on que describe c\'omo una superficie refleja luz seg\'un la direcci\'on de incidencia y de observaci\'on. Es el fundamento matem\'atico del renderizado f\'isicamente basado.

\textit{Draw Call.} Comando emitido por la CPU para que la GPU renderice un lote de geometr\'ia con un material y una pasada concretos. Su reducci\'on es clave para el rendimiento.

\textit{Fragment Shader (\textit{Pixel Shader}).} Programa ejecutado por la GPU para cada fragmento visible, encargado de calcular el color final.

\textit{Frame Time.} Tiempo efectivo requerido para producir un cuadro. En una aproximaci\'on operativa de perfilado suele modelarse como $T_{frame} \approx \max(T_{CPU}, T_{GPU})$, m\'as el posible costo de sincronizaci\'on entre ambos.

\textit{GPU Instancing.} T\'ecnica que permite renderizar varias instancias compatibles de una misma malla y material en menos lotes efectivos. La reducci\'on de \textit{draw calls} depende de las condiciones de compatibilidad y no de una divisi\'on exacta universal.

\textit{Hardware complejo.} Categor\'ia operativa utilizada en esta propuesta para referirse a sistemas electromec\'anicos multicomponente con integraci\'on estructural, energ\'etica, electr\'onica y de control, alta densidad de relaciones espaciales entre piezas y necesidad de inspecci\'on visual detallada del ensamblaje.

\textit{IL2CPP (\textit{Intermediate Language to C++}).} \textit{Backend} de compilaci\'on de Unity que transforma el c\'odigo intermedio (IL) de .NET generado a partir de C\# en C++. En el target Web, ese c\'odigo C++ se compila posteriormente a WebAssembly mediante el \textit{toolchain} correspondiente del build.

\textit{LOD (\textit{Level of Detail}).} Estrategia de optimizaci\'on que intercambia versiones de diferente complejidad geom\'etrica seg\'un distancia o relevancia visual. Se considera una t\'ecnica de referencia del \textit{pipeline}, aunque su adopci\'on final depende del contenido realmente exportado y configurado.

\textit{Metallic Map.} Textura que define qu\'e regiones del material deben tratarse como met\'alicas frente a diel\'ectricas, afectando la reflectancia especular base.

\textit{MVP (Producto M\'inimo Viable).} Versi\'on de un producto con las funcionalidades suficientes para poner a prueba supuestos esenciales con usuarios reales y orientar iteraciones posteriores.

\textit{Normal Map.} Textura que almacena variaciones de normales para simular detalle superficial sin aumentar el conteo de tri\'angulos.

\textit{Occlusion Culling.} T\'ecnica que evita dibujar objetos no visibles porque quedan ocultos desde la perspectiva de la c\'amara. En este proyecto se asume como principio general de optimizaci\'on, no como afirmaci\'on autom\'atica de uso en todos los escenarios de la build final.

\textit{PBR (\textit{Physically-Based Rendering}).} Modelo de iluminaci\'on que aproxima la interacci\'on f\'isica entre luz y materiales mediante conservaci\'on de energ\'ia y teor\'ia de microfacetas.

\textit{Retopolog\'ia.} Proceso de reconstruir la topolog\'ia de una malla para optimizar su estructura geom\'etrica y su viabilidad en tiempo real.

\textit{Roughness Map.} Textura que controla la rugosidad microsc\'opica de la superficie. Valores bajos producen reflejos n\'itidos; valores altos generan reflejos difusos.

\textit{Shader.} Programa ejecutado por la GPU que participa en el c\'alculo visual de v\'ertices o fragmentos.

\textit{Specular.} Componente de reflexi\'on dependiente del \'angulo de observaci\'on y de la microgeometr\'ia del material.

\textit{Texel Density.} Relaci\'on entre la resoluci\'on efectiva de la textura y la longitud real de superficie que representa. Puede expresarse como $\rho = R_{\text{texture}} / L_{\text{mundo}}$ en p\'ixeles por cent\'imetro cuando se toma una direcci\'on o tramo de referencia.

\textit{UV Unwrapping.} Proceso de desplegar una superficie 3D en coordenadas 2D para aplicar texturas con la menor distorsi\'on posible.

\textit{Vertex Shader.} Programa ejecutado por la GPU para cada v\'ertice de la malla, responsable de transformar posiciones entre espacios geom\'etricos.

\textit{WebAssembly (WASM).} Formato binario ejecutable en navegadores que permite desplegar l\'ogica compilada dentro de un modelo de memoria lineal, utilizado por Unity para ejecutar en WebGL el c\'odigo generado a partir de C\#.

\textit{WebGL.} API gr\'afica del navegador basada en OpenGL ES para renderizado 2D y 3D acelerado por hardware sin complementos externos.

\newpage
